\title{Ubiquitin ligases in oncogenic transformation and cancer therapy}
\begin{document}
\maketitle

\begin{abstract}
| The cellular response to external stress signals and DNA damage depends on the activity of ubiquitin ligases (E3s), which regulate numerous cellular processes, including homeostasis, metabolism and cell cycle progression. E3s recognize, interact with and ubiquitylate protein substrates in a temporally and spatially regulated manner. The topology of the ubiquitin chains dictates the fate of the substrates, marking them for recognition and degradation by the proteasome or altering their subcellular localization or assembly into functional complexes. Both genetic and epigenetic alterations account for the deregulation of E3s in cancer. Consequently, the stability and/or activity of E3 substrates are also altered, in some cases leading to downregulation of tumour-suppressor activities and upregulation of oncogenic activities. A better understanding of the mechanisms underlying E3 regulation and function in tumorigenesis is expected to identify novel prognostic markers and to enable the development of the next generation of anticancer therapies. This Review summarizes the oncogenic and tumour-suppressor roles of selected E3s and highlights novel opportunities for therapeutic intervention. NATURE REVIEWS | CANCER VOLUME 18 | FEBRUARY 2018 | 69 REVIEWS © 2 0
\end{abstract}

\section{Recovered Text}
The proteome is the true driver or mediator of cellular
functions and is therefore the preferred target of anti­
cancer therapies. Most targeted therapies are directed
against signalling circuits that are deregulated in cancer,
exemplified by the targeting of protein kinases1. Altered
activity of key regulatory proteins can be due to genetic or
epigenetic modifications, which take place in the course
of cell transformation. Signalling proteins are regulated
by several site-specific post-translational modifications,
of which ubiquitylation is second only to phosphorylation
in abundance2. Ubiquitylation is orchestrated by the
sequential activity of ubiquitin-­activating enzymes (E1s),
ubiquitin-conjugating enzymes (E2s) and ubiquitin ligases
(E3s) (BOX 1). Virtually every cellular protein is subject
to ubiquitylation at least once in its lifetime, exempli­
fying the exquisite homeostatic control of this process.
Ubiquitylation marks proteins for selective protea­
somal or lysosomal degradation or marks entire orga­
nelles for autophagic clearance3–5. The two degradation
pathways show extensive crosstalk and cooperation, as
illustrated by the fact that proteasomes can be cleared
by autophagy and that autophagy substrates can undergo
proteasomal degradation in autophagy-deficient cellu­
lar compartments such as the nucleus6,7. Besides this
fundamental role in maintaining a healthy proteome,
ubiquitylation has multiple non-degradative functions,
including the regulation of protein activity, localization
and complex formation8. Consequently, ubiquitylation is
associated with almost every cellular process, including
the regulation of DNA integrity, gene expression and
metabo­lism8–10. Given the central role of ubiquitylation,
it is not surprising that its deregulation is associated with
a number of diseases, including cancer. The activity of
many E3s is deregulated in cancer (Supplementary infor­
mation S1 (table)) by epigenetic and genetic mechanisms
and/or as a consequence of altered post-translational
mechanisms, which are modified in response to extrinsic
and intrinsic cues (reviewed in REF. 11) (FIG. 1).
Multiple mouse models have been developed to
decipher the role of specific E3s in cancer (TABLE 1).
E3s can elicit oncogenic (for example, MDM2, the E3
for p53) or tumour-suppressive (for example, von
Hippel–Lindau disease tumour suppressor (VHL) and
BRCA1) activity; however, shifting from one function
to the other can be seen upon altered cellular signalling,
as E3s regulate a diverse set of substrates, the regula­
tion of which is influenced by the cellular context (for
example, genetic background, cell lineage, differenti­
ation state and stress level) and subcellular localiza­
tion of the E3 (for example, speckle-type POZ protein
(SPOP), the adaptor for the cullin 3–really interesting
new gene (RING)–E3 ligase (CRL3) complex, functions
as a tumour suppressor in the nucleus but as an onco­
gene in the cytoplasm12,13). Therefore, altered expression
or post-translational modifications play key roles in the
temporal and spatial function of E3s, diverting tumour
suppressor E3s to oncogenes or vice versa. Given the
substrate diversity of E3s, targeting one E3 potentially
1Sanford Burnham Prebys
Medical Discovery Institute,
La Jolla, California 92130,
USA.
2University of Maryland
School of Medicine,
Baltimore, Maryland 21201,
USA.
3Technion Integrated Cancer
Center, Technion, Israel
Institute of Technology
Faculty of Medicine, Haifa
31096, Israel.
Correspondence to Z.A.R.
zeev@ronailab.net
doi:10.1038/nrc.2017.105
Published online 15 Dec 2017
Ubiquitylation
An enzymatic reaction that
leads to the attachment of
ubiquitin moieties either to
ubiquitin itself (creating
polyubiquitin chains) or to
other proteins via isopeptide
linkages.
Ubiquitin-activating
enzymes
(E1s). Enzymes that activate
ubiquitin in an ATP-dependent
manner.
Ubiquitin ligases in oncogenic
transformation and cancer therapy
Daniela Senft1, Jianfei Qi2 and Ze'ev A. Ronai1,3
Abstract | The cellular response to external stress signals and DNA damage depends on the activity
of ubiquitin ligases (E3s), which regulate numerous cellular processes, including homeostasis,
metabolism and cell cycle progression. E3s recognize, interact with and ubiquitylate protein
substrates in a temporally and spatially regulated manner. The topology of the ubiquitin chains
dictates the fate of the substrates, marking them for recognition and degradation by the
proteasome or altering their subcellular localization or assembly into functional complexes. Both
genetic and epigenetic alterations account for the deregulation of E3s in cancer. Consequently, the
stability and/or activity of E3 substrates are also altered, in some cases leading to downregulation
of tumour-suppressor activities and upregulation of oncogenic activities. A better understanding of
the mechanisms underlying E3 regulation and function in tumorigenesis is expected to identify
novel prognostic markers and to enable the development of the next generation of anticancer
therapies. This Review summarizes the oncogenic and tumour-suppressor roles of selected E3s and
highlights novel opportunities for therapeutic intervention.
NATURE REVIEWS | CANCER
 VOLUME 18 | FEBRUARY 2018 | 69
REVIEWS
©
 2
0
1
7
 M
a
c
m
i
l
l
a
n
 P
u
b
l
i
s
h
e
r
s
 L
i
m
i
t
e
d
,
 p
a
r
t
 o
f
 S
p
r
i
n
g
e
r
 N
a
t
u
r
e
.
 A
l
l
 r
i
g
h
t
s
 r
e
s
e
r
v
e
d
.
affects multiple processes required for the malignant
phenotype (FIG. 2), pointing to E3s as desirable drug tar­
gets. However, a deeper, context-dependent understand­
ing of each targeted E3 is required to appreciate possible
tumour-promoting effects of putative inhibitors of E3s.
In this Review, we describe how deregulation of ubiqui­
tylation affects malignant transformation, tumour pro­
gression and therapy resistance. We focus predominantly
on the effects of deregulated ubiquitylation on DNA dam­
age repair, cell cycle regulation, gene expression and signal
Box 1 | The ubiquitin system
The human genome encodes two ubiquitin-activating enzymes (E1s), which
activate ubiquitin (Ub) in an ATP-dependent manner and transfer it to one of
~30 ubiquitin-conjugating enzymes (E2s). The specificity and selectivity of the
ubiquitin-conjugating system are conferred by more than 600 ubiquitin ligases
(E3s), which are sub-classified into three groups according to their mode of
ubiquitin ligation. (a) E3s containing really interesting new gene (RING) and
UFD2 homology (U‑box) domains belong to the first group; the most abundant
(~600 members) is the family of RING E3s. The RING domain coordinates two
Zn2+ ions in a cross-brace arrangement to adopt the structure required for
binding to E2s. RING E3s rely on the enzymatic activity of E2s to ubiquitylate
substrates and can act either independently or as part of multisubunit E3
complexes. Examples of the latter are the cullin–RING–E3 ligase (CRL) family
of complexes, including the S-phase kinase-associated protein 1 (SKP1)–
cullin 1–F‑box protein (SCF) complex and the APC/C (anaphase-promoting
complex; also known as the cyclosome). U‑box E3s (similar to RING E3s)
function as scaffolds for E2s for ubiquitin transfer, but they do not require Zn2+
coordination to adapt their structure. (b) The second group of E3s is the
28‑member homologous to E6AP carboxy terminus (HECT) family, each of
which contains an ~350 amino acid HECT domain that forms a thiol-ester bond
with ubiquitin and then conjugates it to the substrate. According to the
structure of the N‑terminal domain, which serves as the substrate recognition
domain, HECT E3s can be categorized into three subfamilies: NEDD4 and
NEDD4‑like E3s, which contain WW domains; HERC E3s, which harbour
regulator of chromosome condensation 1 (RCC1)‑like domains; and HECT E3s,
which harbour neither WW nor RCC1‑like domains. (c) The third group of
E3s is the 14‑member RING-between-RING (RBR) family. These enzymes have
a RING1–in‑between RING (IBR)–RING2 motif in which the RING1 domain
binds to a ubiquitin-loaded E2 and transfers ubiquitin to the catalytic cysteine
of the RING2 domain, which then conjugates ubiquitin to the substrate. Thus,
RBR E3s function as hybrids of RING E3s and HECT E3s.
The fate of ubiquitylated proteins is largely determined by the ubiquitin
chain topology. In a simplified view, K48‑linked and K11‑linked
ubiquitin chains are usually associated with proteolysis; K63‑linked or
M1‑linked ubiquitin chains mediate the assembly of signalling complexes, as
illustrated by the nuclear factor-κB (NF‑κB) pathway; and monoubiquitylation
serves as a signal for chromatin regulation and protein sorting and trafficking.
Ubiquitin is also phosphorylated, and this modification has been implicated
in the removal of damaged mitochondria by autophagy9,10,188–190. The ubiquitin
code is further modified by the proteolytic activity of ~100 deubiquitylating
enzymes (DUBs) that, for example, remove K48‑chains or K63‑chains to
prevent degradation or abrogate signalling events, respectively10,191. Our
understanding of the role of the individual ubiquitin signals is constantly
evolving. For example, monoubiquitylation was recently suggested to be a
robust signal for proteasomal degradation192.
Decoders of the ubiquitin code (ubiquitin receptors or readers) are
proteins containing at least one of 20 structurally different ubiquitin-
binding domains (UBDs) that serve to recognize linkage-specific
ubiquitin chains and mediate cellular processes (reviewed in REF. 9). As
such, ubiquitin receptors mediate shuttling and binding of ubiquitylated
cargo to the proteasome, initiate autophagy and/or mitophagy and
orchestrate the assembly of signalling complexes9. Many linkage-specific
ubiquitin binding proteins still await identification193.
Ubiquitin-conjugating system
Ubiquitin receptors mediate downstream effects
2 E1s
E1
E2
E2
E2
30 E2s
600 E3s
(a)
Monoubiquitylation
Single
Multiple
DUB
Homotypic
(K6, K11, K27, K29, K33,
K48, K63 and M1)
(e.g. mixed and branched)
Heterotypic
Polyubiquitylation
Ubiquitin code
Ub
Ub
Ub
Ub
Ub
Ub
E3 (HECT)
(b)
E3 (RING, U-box)
Substrate
Signalling complex assembly
Chromatin remodelling
DNA repair
Endocytosis
Autophagy
(e.g. proteins and organelles)
Proteasomal
degradation
Protein conformation and/or activity
Substrate
Substrate
(c)
E3 (RBR)
Ub receptor
REVIEWS
70 | FEBRUARY 2018 | VOLUME 18
www.nature.com/nrc
©
 2
0
1
7
 M
a
c
m
i
l
l
a
n
 P
u
b
l
i
s
h
e
r
s
 L
i
m
i
t
e
d
,
 p
a
r
t
 o
f
 S
p
r
i
n
g
e
r
 N
a
t
u
r
e
.
 A
l
l
 r
i
g
h
t
s
 r
e
s
e
r
v
e
d
.
Ubiquitin-conjugating
enzymes
(E2s). Enzymes that first accept
activated ubiquitin from a
ubiquitin-activating enzyme (E1)
and then transfer it to
substrates.
transduction (FIG. 2), which dictate such pivotal cell fate
decisions as senescence or quiescence, proliferation, differ­
entiation or stemness and cell death. Lastly, given the role
of selected ubiquitin-proteasome pathways in dictating
the malignant phenotype, the opportunities for targeting
deregulated ubiquitylation in cancer are also discussed.
Ubiquitin ligases in genome maintenance
Cell cycle control. Signals that control cell cycle entry, pro­
gression and arrest are commonly deregulated in ­cancer,
and the subsequent disruption of DNA replication,
DNA repair and chromosomal segregation often leads
to genomic instability14. Deregulation of E3s that induce
Figure 1 | Mechanisms underlying deregulated ubiquitylation in cancer.
Deregulated ubiquitylation in cancer can be attributed to epigenetic, genetic,
transcriptional and post-translational mechanisms. Some ubiquitin ligases
(E3s) are encoded by genes that are well recognized to confer susceptibility
for familial cancers, such as the gene encoding von Hippel–Lindau disease
tumour suppressor (VHL) in renal cell carcinoma (RCC)194 or the gene encoding
BRCA1 in breast cancer and ovarian cancer56. Large-scale analyses of cancer
genomes have identified additional E3s that are altered by recurrent
mutations or copy number changes in diverse cancers. The abundance and
activity of E3s can also be regulated by post-translational mechanisms
such as phosphorylation, ubiquitylation or protein–protein interactions,
as demonstrated, for example, with the E3s MDM2 and SIAH2
(REFS 96,98,195,196). The activity and abundance of deubiquitylating enzymes
(DUBs) are also regulated genetically and epigenetically, as recently
reviewed191. In addition to deregulation of E3s and DUBs, the ubiquitin system
is modulated by genetic alterations of the targeted substrates. For example,
recognition of the ubiquitylation sites on MYC143 and the fusion protein
transmembrane protease serine 2 (TMPRSS2)–ETS-related gene (ERG)152
is disrupted by mutations. Ubiquitylation is a dynamic and reversible process
that responds to a variety of internal and external stresses, including DNA
damage and hypoxic, oxidative and metabolic stresses, which are all
encountered by cancer cells during malignant transformation, during
metastatic dissemination and in response to therapy. Each individual alteration
in the ubiquitin system can have a profound effect on the regulation of
cancer-associated pathways by modulating the localization, activity, signalling
complex formation and abundance of major regulatory hubs. AKR1C3,
aldo-keto reductase family 1 member C3; APC/C, anaphase-promoting
complex; also known as the cyclosome; BAP1, BRCA1‑associated protein 1;
BTRC, encoding β-TRCP; CYLD, cylindromatosis; E2, ubiquitin-conjugating
enzyme; FBXO11, F-box only protein 11; FBXW7, F‑box/WD repeat-containing
protein 7; GSK3, glycogen synthase kinase 3; HECT, homologous to E6AP
carboxy terminus; KEAP1, kelch-like ECH-associated protein 1; NRF2, nuclear
factor erythroid 2‑related factor 2; PARK2, encoding parkin; RBR,
RING-between-RING; RING, really interesting new gene; SIAH, seven in
absentia homolgue; SKP2, S‑phase kinase-associated protein 2; SPOP,
speckle-type POZ protein; STUB1, STIP1 homology and U box-containing
protein 1 (also known as CHIP); TRIM7, tripartite motif 7; U-box, UFD2
homology; Ub, ubiquitin; USP, ubiquitin carboxyl-terminal hydrolase.
DUB
Mechanisms of deregulation of ubiquitylation in cancer
Consequences of aberrant ubiquitylation in cancer
Abundance, function and
localization of E3s
Substrate
Properties of the substrate
Activity of DUBs
Aberrant ubiquitin code
Substrate localization
Regulation of hallmark cancer processes
Substrate interactions
Substrate activity
Substrate abundance
E3
(RING, U-box, HECT
and RBR domains)
monomer, dimer, complex
Genetic
(Post-)transcriptional
Post-translational
Coding mutation
  SPOP, FBXW7, VHL, KEAP1,
  BRCA1, APC/C, HUWE1
Non-coding mutation
  MDM2, BTRC
Amplification
  MDM2, SKP2, NEDD4
Deletion
  FBXW7, VHL, PARK2,
  BRCA1, HUWE1, FBXO11
Genetic
Mutation
  MYC, JUN, NRF2
Fusion
  TMPRSS2–ERG
Genetic
Deletion
  BAP1, CYLD
Amplification
  USP15
Promoter hypermethylation
  VHL, KEAP1, BRCA1, Parkin,
  STUB1
Transcription factor activity
  SKP2, SIAH1, SIAH2, RNF125,
  SPOP
miRNAs
  FBXW7, mutant BTRC, RNF8
Splicing
  MDM2, FBXW7
Post-translational
Kinase activity
  GSK3, BRAF
Phosphorylation
  MDM2, TRIM7, SIAH2
(Self-)ubiquitylation
  SIAH2, MDM2, NEDD4
Protein–protein interaction
  MDM2–MDM4, SIAH–AKR1C3
Neddylation, acetylation,
glycosylation
E2
Ub
P
Post-translational
Phosphorylation
  USP20, USP4
Protein–protein interaction
  BAP1
(Post-)transcriptional
USP13, USP22
REVIEWS
NATURE REVIEWS | CANCER
 VOLUME 18 | FEBRUARY 2018 | 71
©
 2
0
1
7
 M
a
c
m
i
l
l
a
n
 P
u
b
l
i
s
h
e
r
s
 L
i
m
i
t
e
d
,
 p
a
r
t
 o
f
 S
p
r
i
n
g
e
r
 N
a
t
u
r
e
.
 A
l
l
 r
i
g
h
t
s
 r
e
s
e
r
v
e
d
.
Ubiquitin ligases
(E3s). Enzymes that facilitate
substrate recognition and
guide the transfer of activated
ubiquitin from the ubiquitin-
conjugating enzyme (E2) to
specific substrates.
Mitophagy
The selective autophagic
clearance of damaged
mitochondria.
Table 1 | Examples of ubiquitin ligases in transgenic mouse models
E3s
Phenotypes in transgenic mouse models
Refs
β-TRCP
Btrc overexpression in the mouse mammary gland increased cell proliferation and induced
breast tumour development, which was associated with activation of the NF‑κB pathway
201
Overexpression of a dominant-negative mutant of Btrc (in which the F-box domain is deleted)
in transgenic mice resulted in the development of intestinal, hepatic or urothelial tumours,
which were associated with accumulation of β‑catenin
202
BRCA1
Conditional knockout of Brca1 in mouse mammary epithelial cells resulted in the development
of basal-like triple-negative breast carcinomas
203
Co‑deletion of Brca1 and Trp53 in mouse mammary epithelial cells accelerated the formation
of breast carcinomas
204,205
CDC20
Cdc20‑knockout mice treated with a two-stage carcinogenesis protocol showed a massive
arrest in metaphase and apoptosis of skin tumour cells
19
CDH1
Aged Cdh1‑heterozygous mice show increased susceptibility to spontaneous epithelial
tumours in various organs (for example, mammary gland, lung, liver, kidney, testis and
sebaceous gland)
15
FBXW7
Fbxw7‑knockout mice showed a higher incidence of γ‑irradiation-induced tumour formation
than wild-type mice. Loss of Fbxw7 also changed the spectrum of tumours that developed in
irradiated Trp53‑deficient mice to include epithelial tumours from the lung, liver and ovary
26
Conditional T cell lineage-specific knockout of Fbxw7 in mice resulted in thymic lymphoma
and acute leukaemia
164
Knockout of Fbxw7 in ApcMin/+ mice accelerated intestinal tumour development
206,207
Double knockout of Fbxw7 and Trp53 in the mouse intestine caused aggressive and metastatic
adenocarcinomas that resembled human advanced colorectal cancer
24
HUWE1
Knockout of Huwe1 accelerated tumour formation in a two-stage skin carcinogenesis model,
which was reversed by concomitant knockout of Myc
138
Knockout of Huwe1 in ApcMin/+ mice accelerated intestinal tumour development, which was
associated with increased MYC levels, rapid accumulation of DNA damage and loss of the
second copy of Apc
140
MDM2
Overexpression of Mdm2 in mice induced spontaneous tumour formation. The frequencies of
tumour formation in Mdm2‑transgenic and Trp53‑null mice were statistically indistinguishable
from each other. Mdm2‑transgenic mice (regardless of their Trp53 status) showed a higher
percentage of sarcomas than Trp53‑null mice
208
Parkin
Park2‑knockout mice showed increased hepatocyte proliferation and developed macroscopic
hepatic tumours resembling hepatocellular carcinoma
209
Park2‑knockout mice accelerated lymphoma development upon γ‑irradiation
210
Knockout of Park2 in ApcMin/+ mice increased the development of intestinal tumours
50
SIAH2
Knockout of Siah2 inhibited the development of prostate neuroendocrine tumours and
promoted the regression of atypical hyperplasia upon castration in the TRAMP model
102,154
Knockout of Siah2 in MMTV–PyMT mice delayed the progression of mammary tumours and
was associated with vascular normalization in the tumour microenvironment
104
SKP2
Transgenic overexpression of Skp2 in the mouse prostate induces hyperplasia, dysplasia and
low-grade carcinoma
31
Co‑expression of Skp2 with NrasG12V or myristoylated Akt1 in the mouse liver resulted in
hepatocellular carcinoma
32
Skp2‑knockout inhibited tumour development in a Pten–/–Trp53–/– mouse prostate cancer
model via activation of p27KIP1-dependent, p21CIP1-dependent and ATF4‑dependent
senescence
33
Co‑deletion of Skp2, Rb1 and Trp53 blocked tumorigenesis in the mouse pituitary and prostate
in a p27KIP1‑dependent manner
34
Skp2‑knockout inhibited DMBA–TPA-induced skin tumorigenesis in mice
211
SPOP
Prostate-specific knockout of Spop resulted in the upregulation of Myc and promoted
development of prostatic intraepithelial neoplasia
141
APC, adenomatous polyposis coli; ATF4, activating transcription factor 4; Btrc, encoding β-TRCP; CDC20, cell division cycle 20;
CDH1, CDC20-like protein 1 (also known as FZR1); DMBA, dimethylbenzanthracene; E3s, ubiquitin ligases; FBXW7, F‑box/WD
repeat-containing protein 7; MMTV–PyMT, mouse mammary tumour virus–polyoma middle T antigen; NF‑κB, nuclear factor-κB;
Parkin, encoded by Park2; SIAH2, seven in absentia homologue 2; SKP2, S‑phase kinase-associated protein 2; SPOP, speckle-type
POZ protein; TPA, 12‑O-tetradecanoylphorbol‑13‑acetate; TRAMP, transgenic adenocarcinoma of the mouse prostate.
REVIEWS
72 | FEBRUARY 2018 | VOLUME 18
www.nature.com/nrc
©
 2
0
1
7
 M
a
c
m
i
l
l
a
n
 P
u
b
l
i
s
h
e
r
s
 L
i
m
i
t
e
d
,
 p
a
r
t
 o
f
 S
p
r
i
n
g
e
r
 N
a
t
u
r
e
.
 A
l
l
 r
i
g
h
t
s
 r
e
s
e
r
v
e
d
.
Cyclins
A family of regulatory proteins
that show oscillating
expression throughout the cell
cycle and that are required
for the activation of
cyclin-dependent kinases.
Cyclin-dependent kinase
(CDK). A type of protein that
belongs to a group of serine
and threonine kinases that
require cyclins for activation
and regulate cell cycle
progression.
proteasomal degradation of cyclins and cyclin-­dependent
kinase (CDK) inhibitor proteins and/or that regulate the
assembly of the DNA damage repair machinery thus
contributes to the sustained proliferation and genomic
instability observed in cancer cells (FIG. 3). Among the
best-studied E3s that regulate cell cycle progression are
the APC/C (anaphase-promoting complex; also known
as the cyclosome) and its ­co‑activators cell division
cycle 20 (CDC20) or CDC20-like protein 1 (CDH1; also
known as FZR1) and S-phase kinase-­associated protein 1
(SKP1)–cullin 1–F‑box protein (SCF) (and its component
F‑box proteins F‑box/WD repeat-containing protein 7
(FBXW7), SKP2 and ­β-transducin repeat-containing
protein (β-TRCP; also known as BTRC)) complexes.
E3s deregulated
in cancer
• Receptor recycling
• Protein recruitment
• Complex formation
• Signal transduction
• SCF–FBXW7
• NEDD4
• WWP1
• CBL
• SPOP
• SIAH1 and
SIAH2
• SCF–SKP2
• TRIM28
• Parkin
• STUB1
• RNF5
• SIAH1 and
SIAH2
• HUWE1
• Parkin
• SCF–β-TRCP
• SCF–FBXW7
• CRL3–KEAP1
• SIAH1 and
SIAH2
• STUB1
• SCF–β-TRCP
• HRD1
• GP78
• SCF–FBXW7
• SCF–BTRC
• TRIM28
• MDM2
• Parkin
• BRCA1
• HUWE1
• SCF–SKP2
• MDM2
• APC/C
• SCF–FBXW7
• Cell cycle regulation
• Epigenetic marks
• Transcriptional control
• Genome integrity
• Mitophagy
• Cell death
• Mitochondrial
   biogenesis and
   dynamics
• Redox control
• Proteasomal
   degradation
• Proteostasis
Autophagy
Autophagosome
Mitochondrion
Translation
ER
Most E3s
• ER homeostasis
• Protein secretion
Nucleus
Figure 2 | Cellular processes affected by deregulated ubiquitylation in cancer. Representative ubiquitin ligases (E3s)
that are deregulated in cancer and the biological processes expected to be affected are depicted. As E3s ubiquitylate a
diverse set of substrates, E3 loss or gain of function affects multiple cellular processes simultaneously. For example, S-phase
kinase-associated protein 1 (SKP1)–cullin 1–F‑box protein (SCF)–F‑box/WD repeat-containing protein 7 (FBXW7) targets
cell cycle regulators (for example, cyclin E), oncogenic transcription factors (for example, MYC), cell surface receptors
(for example, NOTCH1), signalling molecules (for example, mTOR) and apoptosis regulators (for example, myeloid cell
leukaemia 1 (MCL1)) for proteasomal degradation. Therefore, SCF–FBXW7 loss of activity through mutations or deletions
leads to genomic instability, increased proliferation and survival and the rewiring of transcriptional and signalling
programmes that affect cancer cell migration, metabolism and stemness. Similarly, SCF–β-transducin repeat-containing
protein (β-TRCP) can serve as a signalling hub to coordinate increased protein synthesis and pro-survival signals following
pro-growth stimuli in normal cells and in cancer cells197–199. Thus, the β-TRCP signalling circuits provide a platform for
therapeutic intervention, for example, in cancers characterized by activation of mTOR signalling. APC/C, anaphase-
promoting complex; also known as the cyclosome; CRL3, cullin 3–really interesting new gene (RING)–E3 ligase;
ER, endoplasmic reticulum; HRD1, also known as synoviolin; GP78, also known as AMFR; KEAP1, kelch-like ECH-associated
protein 1; SIAH, seven in absentia homologue; SKP2, S‑phase kinase-associated protein 2; SPOP, speckle-type POZ protein;
STUB1, STIP1 homology and U box-containing protein 1; TRIM28, tripartite motif 28.
REVIEWS
NATURE REVIEWS | CANCER
 VOLUME 18 | FEBRUARY 2018 | 73
©
 2
0
1
7
 M
a
c
m
i
l
l
a
n
 P
u
b
l
i
s
h
e
r
s
 L
i
m
i
t
e
d
,
 p
a
r
t
 o
f
 S
p
r
i
n
g
e
r
 N
a
t
u
r
e
.
 A
l
l
 r
i
g
h
t
s
 r
e
s
e
r
v
e
d
.
a
b
• Functions as tumour suppressor
• Downregulation mediates
   premature S-phase entry and
   genomic instability
• Deleted or inactivated in cancer
• Premature S-phase entry
• DNA-repair defects
• Overexpressed in cancer
• p27KIP1-dependent and p27KIP1-
   independent tumour promoter
• Inhibits terminal differentiation
• Deleted or inactivated in cancer
• Associated with enforced cell
   cycle progression and genomic
   instability
DSBs
DSB repair
NHEJ
HR
• G1–S checkpoint
• S checkpoint
• G2–M checkpoint
Suppressed during G1 by
Ub-dependent mechanisms
Major repair pathway in
HR-deficient cancers;
error prone; source of
genomic instability
Inactivated in cancer owing to
germline or somatic mutations in
BRCA1 or BRCA2 or genes encoding
other HR-regulatory proteins
Restriction point
ATM and ATR
CHK1
CHK1
and CHK2
USP11
CDC25A,
CDC25B and
CDC25C
MDM2
p53
Crosstalk with DNA damage response and
DNA repair machinery
• Role in cancer unclear
• Tumour-promoting or
   tumour-suppressing
   activities
Therapeutic target to
prevent mitotic exit
• Cyclin E
• Aurora kinase A
• Cyclin D
• Cyclin E
• Cyclin B1
Restriction
point
Cyclin D, CDK4
and CDK6
Cyclin E,
CDK2
Cyclin A, CDK1
and CDK2
Cyclin B and CDK1
M
G1
G2
S
SCF–FBXW7
• WEE1
• EMI1
• Claspin
• MDM2
• CDC25A
• CDC25B
• REST
• CEP68
• eEF2K
• FANCM
• CDC20
• Aurora kinases
• CDC25A
• PLK1
• CTIP
• Cyclin A
• Cyclin B1
• Securin
• p21CIP1
• Cyclins
• Spindle proteins
• SKP2
• Claspin
SCF–β-TRCP
SCF–β-TRCP
APC/C–CDH1
CRL3–KEAP1
APC/C–CDH1
APC/C–CDC20
• p27KIP1
• p21CIP1
SCF–SKP2
Parkin
G1–S
checkpoint
G2–M checkpoint
SAC
S checkpoint
p21CIP1
CDK
CDK
REVIEWS
74 | FEBRUARY 2018 | VOLUME 18
www.nature.com/nrc
©
 2
0
1
7
 M
a
c
m
i
l
l
a
n
 P
u
b
l
i
s
h
e
r
s
 L
i
m
i
t
e
d
,
 p
a
r
t
 o
f
 S
p
r
i
n
g
e
r
 N
a
t
u
r
e
.
 A
l
l
 r
i
g
h
t
s
 r
e
s
e
r
v
e
d
.
Haploinsufficient
A state in which one copy of a
gene is inactivated or deleted
and the remaining functional
copy is not sufficient to
preserve normal function.
Replication origins
Sites in the DNA where the
replication machinery is loaded
at the onset of DNA synthesis.
Spindle assembly
checkpoint
(SAC). A stage in the cell cycle
that is activated during mitosis
and meiosis to delay cell
division until all chromosomes
are correctly attached to the
spindle.
Phosphodegron
One or multiple
phosphorylated residues in a
protein substrate that are
necessary for recognition by
some ubiquitin ligases.
Although APC/C itself is rarely mutated in cancer,
increasing evidence supports a tumour-­suppressor
role for CDH1 and an oncogenic role for CDC20
(Supplementary information S1 (table)). CDH1
functions as a haploinsufficient tumour suppressor
(FIG. 3a). Aged Cdh1+/–, but not wild-type (WT), mice
show increased susceptibility to spontaneous epi­
thelial tumours in various organs15. Accordingly,
Cdh1‑deficient mouse embryonic fibroblasts (MEFs)
show reduced proliferation and accumulate chromo­
somal aberrations15. CDH1 knockdown in human
bone osteosarcoma U2OS cells causes accumulation
of cyclin A and cyclin B, premature entry into S‑phase
with reduced loading of pre-replication complexes
onto DNA replication origins and accumulation of DNA
­double strand breaks (DSBs) during mitosis owing to
the presence of replication intermediates16.
Interest in the oncogenic potential of CDC20 arose
when it was discovered that residual CDC20 activ­
ity upon activation of the spindle assembly checkpoint
(SAC), (which sequesters CDC20 into the mitotic
checkpoint complex, consisting of BUBR1 (also known
as BUB1B), BUB3 and MAD2, to delay cyclin B1 degra­
dation and induce mitotic arrest) promotes escape from
antimitotic drug-induced apoptosis17,18. Consequently,
blocking mitotic exit via CDC20 inhibition emerged as
a more efficient means to induce apoptosis than spindle
checkpoint-­dependent mitotic poisons18. In a two-stage
skin carcinogenesis mouse model, localized deletion of
Cdc20 results in massive metaphase arrest and apop­
tosis, and this phenotype is also observed in vitro in
Cdc20–/– MEFs transformed with oncogenic RASG12V
and early region 1A (E1A) of human adenovirus type 5
(REF. 19). These findings prompted the development of
two APC/C–CDC20 inhibitors: the small-molecule
inhibitor, tosyl‑l‑arginine methyl ester (TAME), which
binds to APC/C and prevents its activation by CDC20 or
CDH1 (REF. 20), and apcin, which binds to CDC20 and
blocks its interaction with the destruction box (D‑box),
which is present in almost all APC/C substrates21; apcin
can act synergistically with TAME to block mitotic
exit, exemplifying the difficulties in pharmacologi­
cally blocking the activity of the multisubunit APC/C
efficiently21. Given that both inhibitors block CDH1 as
well as CDC20 (REF. 20), caution is required to monitor
for possible tumour-promoting consequences of CDH1
inhibition. The development of novel CDC20‑specific
inhibitors could eliminate this concern. Recent advances
in understanding how phosphorylation of the APC/C
subunits stimulates CDC20 loading and activation of
APC/C are an important step in this direction22,23.
F‑box proteins are often deregulated in cancer
(Supplementary information S1 (table)), thereby affect­
ing SCF complex activity (FIG. 3). FBXW7‑containing SCF
complexes mediate degradation of cyclin E (FIG. 3a), and
impaired SCF–FBXW7 function leads to sustained pro­
liferation and genomic instability24,25. FBXW7 functions
as a p53‑dependent haploinsufficient tumour suppres­
sor; thus, the effects of FBXW7 substrate accumula­
tion upon heterozygous inactivation of FBXW7 can be
reversed by p53 expression24–26. In a mouse model, intes­
tinal co‑deletion of Fbxw7 and Trp53 results in advanced
adenocarcinomas with increased cyclin E expression and
a phenotype of chromosomal instability24. In addition,
radiation-induced lymphomas in Trp53+/– mice but not
Trp53–/– mice showed loss of heterozygosity and a 10\%
mutation rate of Fbxw7 (REF. 26). However, rather than
leading to cyclin E accumulation, deregulation of Aurora
kinase A, another substrate of SCF–FBXW7, mediated
genomic instability in this context26.
The oncogenic SCF–SKP2 complex regulates a
number of CDK inhibitors, of which the best studied
is the tumour suppressor p27KIP1 (REFS. 27–30) (FIG. 3a).
Transgenic overexpression of Skp2 in the mouse pros­
tate induces hyperplasia, dysplasia and low-grade car­
cinoma31, while co‑expression of Skp2 with NrasG12V or
myristoylated Akt1 in the mouse liver results in hepato­
cellular carcinoma (HCC)32. Conversely, Skp2‑knockout
efficiently inhibits tumour development in a condi­
tional Pten-deficient and Trp53‑deficient mouse pros­
tate cancer model via activation of p27KIP1-dependent,
p21CIP1-dependent and activating transcription factor 4
(ATF4)-dependent senescence33. Co‑deletion of Skp2,
Rb1 and Trp53 blocks tumorigenesis in the mouse
­pituitary and prostate in a p27KIP1‑dependent manner34.
Figure 3 | Ubiquitin ligases coordinate the cell cycle and DNA damage repair to
maintain genome integrity. a | Phosphorylation and ubiquitylation coordinate the
temporal activity of cyclin-dependent kinase (CDK)–cyclin complexes, thereby mediating
cell cycle progression and checkpoint control. Displayed are some of the ubiquitin ligases
(E3s), which are well-known cell cycle regulators, such as the APC/C (anaphase-promoting
complex; also known as the cyclosome) and S-phase kinase-associated protein 1 (SKP1)–
cullin 1–F‑box protein (SCF)–F-box/WD repeat-containing protein 7 (FBXW7) complexes
and parkin, whose function in the cell cycle is emerging. APC/C primarily mediates
progression through mitosis by temporally coordinating the recruitment of co‑activators
(cell division cycle 20 (CDC20) or CDC20-like protein 1 (CDH1)). Thus, APC/C–CDC20
mediates anaphase entry, while APC/C–CDH1 acts during mitotic exit and early G1. APC/C
is also regulated by phosphorylation, binding of inhibitory molecules such as early mitotic
inhibitor 1 (EMI1; also known as FBXO5) and the spindle assembly checkpoint (SAC)200. SCF
is a four-protein complex consisting of the scaffold cullin 1, a really interesting new gene
(RING) domain-containing component RBX1, the SKP1 adaptor protein and one of ~68
various F‑box proteins (for example, FBXW7, SKP2 or β-transducin repeat-containing
protein (β‑TRCP)). Most F‑box proteins mediate the recognition of substrates by binding
to their phosphodegron motifs (although phosphorylation-independent mechanisms of
substrate recognition exist), which contain residues that can be phosphorylated by
multiple kinases. The SCF complex is active throughout the cell cycle; SCF substrates
include a subset of cyclins and CDK inhibitors, through which the SCF complex regulates
progression from G1 to the onset of mitosis. Some example substrates of each E3 are
displayed in the cream boxes, and the effects of altered E3 function in cancer are indicated
in the grey boxes. b | DNA double strand breaks (DSBs) mediate activation of DNA damage
sensors ataxia telangiectasia mutated (ATM), ataxia telangiectasia and Rad3‑related (ATR),
checkpoint kinase 1 (CHK1) and CHK2. These in turn inactivate the E3 MDM2, leading to
p53 stabilization and/or promoting SCF–β-TRCP-mediated degradation of CDK
phosphatases CDC25A, CDC25B and CDC25C; both pathways converge on the
attenuation of CDK activity. ATM and ATR also initiate the recruitment of the DNA repair
machinery to the sites of DNA damage, which is in turn under the control of ubiquitylation.
Suppression of homologous recombination (HR) during G1 involves competition between
p53‑binding protein 1 (53BP1) and BRCA1 at DSBs, APC/C–CDH1‑mediated CtBP-
interacting protein (CTIP) degradation and cullin 3–RING–E3 ligase (CRL3)–kelch-like
ECH-associated protein 1 (KEAP1)-dependent inhibition of BRCA1–partner and localizer
of BRCA2 (PALB2)–BRCA2 complex formation62–64. CEP68, centrosomal protein of 68 kDa;
eEF2K, eukaryotic elongation factor 2 kinase; FANCM, Fanconi anaemia group M protein;
NHEJ, non-homologous end joining; PLK1, polo-like kinase 1; REST, repressor-element 1
(REI)-silencing transcription factor; Ub, ubiquitin; USP11, ubiquitin carboxyl-terminal
hydrolase 11.
◀
REVIEWS
NATURE REVIEWS | CANCER
 VOLUME 18 | FEBRUARY 2018 | 75
©
 2
0
1
7
 M
a
c
m
i
l
l
a
n
 P
u
b
l
i
s
h
e
r
s
 L
i
m
i
t
e
d
,
 p
a
r
t
 o
f
 S
p
r
i
n
g
e
r
 N
a
t
u
r
e
.
 A
l
l
 r
i
g
h
t
s
 r
e
s
e
r
v
e
d
.
Destruction box
(D‑box). A conserved sequence
of amino acids (RxxL) in
proteins that is recognized by
the APC/C (anaphase-
promoting complex; also
known as the cyclosome).
Micronuclei
Extranuclear bodies that form
if chromosome fragments or
entire chromosomes are not
incorporated into the nucleus
following cell division.
β-TRCP-containing SCF complexes play dual roles
in cell cycle checkpoint control: they mediate cell cycle
arrest via degradation of checkpoint kinase 1 (CHK1)-
phosphorylated CDC25A35 and relieve the arrest via
degradation of WEE1, claspin, eukaryotic elongation
factor 2 kinase (eEF2K) and Fanconi anaemia group M
protein (FANCM) following phosphorylation by major
M‑phase kinases, such as polo-like kinase 1 (PLK1) and
CDK1 (REFS 36–39) (FIG. 3a). In addition, SCF–β-TRCP
promotes cell cycle arrest by targeting the degrada­
tion of casein kinase I (CKI)-phosphorylated MDM2,
which leads to stabilization of p53 (REF. 40). Although
these mechanisms might suggest a tumour-suppres­
sor function for SCF–β-TRCP, the situation in human
cancers is not clear. In part, this may be explained by
its context-dependent opposing functions in cell cycle
progression and arrest. However, many SCF–β-TRCP
substrates are themselves tumour suppressors (for exam­
ple, inhibitor of nuclear factor‑κB (IκB), forkhead box
protein O3 (FOXO3), p19ARF and repressor-element 1
(REI)-silencing transcription factor (REST)) and onco­
genes (for example, TWIST1, MDM2 and β‑catenin)40–46,
suggesting that the role of SCF–β-TRCP in cell cycle
control may not entirely explain its effects in cancer.
In addition to the classical cell cycle regulators,
­studies in genetic mouse models and analysis of muta­
tions in human cancers have pointed to the role of
the E3 parkin in cell cycle control47–52 (FIG. 3a). Genetic
alteration (including mutations and copy number loss)
of PARK2 (the gene encoding parkin; also known as
PRKN) suggests a tumour-suppressive function for
this E347,48,51. In a mouse model of intestinal tumori­
genesis, the rate of development of intestinal adeno­
mas was found to be approximately fourfold higher in
Park2+/–ApcMin/+ mice than in ApcMin/+ mice; increased
tumour formation was associated with the loss of the
inhibitory effect of parkin on cell proliferation50. In
addition, knockdown of parkin in cancer cell lines is
associated with multi­polar spindles and the formation
of ­micronuclei owing to cyclin E accumulation51. Notably,
analysis of pan-­cancer mutation data revealed that muta­
tions in PARK2 are mutually exclusive with alterations
in G1 phase and S phase cell cycle regulators (cyclin D,
cyclin E and CDK4)48. Furthermore, parkin can regu­
late cyclin D and cyclin E via formation of F-box only
protein 4 (FBXO4)-containing or FBXW7‑containing
parkin–cullin–RING complexes, respectively48. Finally,
parkin may interfere with cell cycle control by associat­
ing with CDC20 or CDH1, promoting degradation of
several key mitotic regulators, including PLK1, NEK2,
cyclin B1, securin, Aurora kinase A and Aurora kinase
B, in an APC/C‑independent manner49. In summary, the
accumulation of mitotic regulators upon parkin loss of
function may contribute to genomic instability, thereby
promoting tumour formation.
In contrast to cancer, where parkin defects are associ­
ated with increased proliferation and increased tumori­
genesis, in Parkinson disease, loss-of-function mutations
of PARK2 impair mitophagy, resulting in accumula­
tion of damaged mitochondria and induction of apop­
tosis53. This difference might be explained in part by the
post-mitotic nature of neurons compared with the pro­
liferative phenotype of genetically unstable cancer cells,
which may be able to circumvent apoptosis and tolerate
mitochondrial damage; however, the roles of mitophagy
and mitochondrial integrity in tumours, in which parkin
function is lost through mutations, remain to be clarified
(reviewed in REF. 54).
DNA damage repair. Among the E3s, MDM2 and
BRCA1 are known to link regulation of the DNA
damage response and cell cycle checkpoints to cancer
development (reviewed in REFS 55,56) (FIG. 3b). Briefly,
MDM2 is overexpressed in a variety of cancers and
promotes tumorigenesis primarily by targeting the
degradation of p53, although the regulation of other
substrates by MDM2 may contribute55 (Supplementary
information S1 (table)). Inhibitors that disrupt the inter­
action between p53 and MDM2 and/or the homologue
MDMX (also known as MDM4) were developed57,
including the cis-imidazoline analogues (otherwise
known as the nutlins) such as Nutlin‑3 and RG7112,
which are currently being assessed in clinical trials for
haematological malignancies58.
BRCA1 forms a heterodimer with BRCA1‑associated
RING domain protein 1 (BARD1) and mediates mono­
ubiquitylation or non-degradative polyubiquitylation
of its substrates56. The BRCA1–BARD1 complex is
implicated in multiple cellular processes by virtue of
its broad range of substrates, including histones, CtBP-
interacting protein (CTIP; also known as RBBP8),
oestrogen ­receptor‑α (ERα), RNA polymerase II
(RNAPII) and transcription initiation factor IIE (TFIIE)
(Supplementary information S1 (table)). The aberrant
role of BRCA1–BARD1 in homologous recombination
(HR) and cell cycle control, and the resulting genomic
instability, is considered to be the key determinant in the
aetiology of breast and ovarian cancer in women carrying
BRCA1 or BRCA2 mutations56. Nonetheless, the impor­
tance of the E3 activity of BRCA1 in these processes is
controversial. For example, a BRCA1–RING-domain-
mutant that lacks ligase activity cannot restore a cell cycle
checkpoint or reverse γ‑irradiation hypersensitivity in
BRCA1‑null human breast cancer cell lines59. Similarly,
mice harbouring the clinically relevant missense Brca1C61G
mutation, which confers aberrant E3 activity and reduces
the interaction with BARD1, exhibit genomic instabil­
ity and tumour development similar to that in Brca1–/–
mice60. However, mice harbouring a Brca1I26A missense
mutation (not found in human tumours), which allows
heterodimer formation with BARD1 but disrupts the
E3 activity of BRCA1, do show reduced tumour devel­
opment compared with that in WT mice61, suggesting
that the E3 activity is not essential for tumour suppres­
sion. As BRCA1 functions as a scaffold protein for mul­
tiple protein complexes, protein–protein interactions
and E3 activity are both likely to play ­important roles in
BRCA1‑mediated tumour suppression.
BRCA1 further links DNA DSB repair to the cell
cycle56,62. DSB repair is regulated by ubiquitylation
and deubiquitylation at multiple levels, starting with
the recruitment of E3s RNF8 and RNF138 to DSBs,
REVIEWS
76 | FEBRUARY 2018 | VOLUME 18
www.nature.com/nrc
©
 2
0
1
7
 M
a
c
m
i
l
l
a
n
 P
u
b
l
i
s
h
e
r
s
 L
i
m
i
t
e
d
,
 p
a
r
t
 o
f
 S
p
r
i
n
g
e
r
 N
a
t
u
r
e
.
 A
l
l
 r
i
g
h
t
s
 r
e
s
e
r
v
e
d
.
Non-homologous end
joining
(NHEJ). An error-prone
mechanism to repair DNA
double strand breaks whereby
the broken ends can be ligated,
even with little or no sequence
complementarity.
Base excision repair
(BER). A DNA repair
mechanism that replaces bases
that are damaged as a result of
oxidation, deamination or
alkylation.
Deubiquitylating enzyme
(DUB). A protease (cysteine
protease or metalloproteinase)
that cleaves the isopeptide
linkage between the protein
substrate (which can be
ubiquitin itself) and the
ubiquitin residue.
followed by a wave of ubiquitin signals that promote the
recruitment of DNA repair factors to mediate chroma­
tin remodelling at DSB sites62. Interestingly, the mode of
DNA repair (that is, non-homologous end ­joining (NHEJ) or
HR) is under cell cycle control (FIG. 3b). HR is restricted
to the cell cycle phases when a sister chromatid is avail­
able for recombination and is thus suppressed during G1,
partially via ubiquitin-dependent mechanisms62–64 (FIG. 3).
These mechanisms might be of therapeutic importance
in cancer, because defective HR renders cells suscepti­
ble to inhibition of base excision repair (BER) mediated
by poly (ADP-ribose) polymerase 1 (PARP1)65–67: the
­deubiquitylating enzyme (DUB) ubiquitin carboxyl-­
terminal hydrolase 11 (USP11) deubiquitylates partner
and localizer of BRCA2 (PALB2) during S and G2 phases
following DNA damage, allowing the formation of the
BRCA1–PALB2–BRCA2 complex and HR repair to
advance in these phases of the cell cycle64. USP11 is often
overexpressed in cancer, confers resistance to PARP inhib­
itors68 and may serve as a biomarker for PARP-inhibitor
resistance. Conversely, its targeting may ­sensitize resistant
tumours to PARP inhibition68.
Furthermore, SCF–FBXW7 has been suggested to
play a direct role in DNA DSB repair69. Activation of
ataxia telangiectasia mutated (ATM), a crucial medi­
ator of the DNA damage response, leads to phosphory­
lation of SCF–FBXW7 and its recruitment to the DSB
sites, followed by K63‑linked polyubiquitylation of X‑ray
repair cross-complementing protein 4 (XRCC4), a repair
protein implicated in NHEJ. K63‑ubiquitylated XRCC4
enhances its association with the KU70 (also known as
XRCC6) and KU80 (also known as XRCC5) complex,
thereby increasing NHEJ repair69.
These examples establish the role of deregulated E3s
in the uncontrolled proliferation and genomic instability
that drive malignant transformation, tumour progres­
sion and therapy resistance. Individual E3s have a broad
spectrum of substrates and thus can play multifaceted
roles by serving as a nexus to coordinate cell growth,
proliferation and survival under both favourable and
hostile growth conditions (FIG. 2).
Signal transduction regulation by ubiquitin ligases
E3s can regulate major growth-promoting pathways,
including those targeted by current anticancer therapies
such as the MAPK or PI3K–AKT–mTOR pathways1,8
(FIG. 4a). Sustained activation of pathways that promote
growth and survival constitutes a stressful environment;
consequently, cancer cells must ensure that metabolic
processes and stress signalling pathways are coordi­
nately regulated in order to overcome these potentially
­deleterious conditions.
Regulation of MAPK signalling. The MAPK and PI3K–
AKT–mTOR pathways cooperate to drive cell growth,
proliferation and survival70 (FIG. 4a). They are the most
hyperactivated pathways in cancer, and, not surprisingly,
they are tightly regulated by the ubiquitin machinery.
For example, the abundance of tyrosine kinase receptors
at the cell surface is regulated by ubiquitin-dependent
recycling (reviewed in REF. 4). Degradation of RAS by the
E3 NEDD4 is part of a negative feedback loop in which
RAS signalling induces transcriptional upregulation of
NEDD4 that in turn limits the activity of WT RAS but
not oncogenic RAS71. As NEDD4 is a known negative
regulator of PTEN72, NEDD4 upregulation increases
PTEN degradation, further enhancing the malignant
phenotype of RAS-driven tumours71. The RAS–NEDD4
relationship exemplifies how the genetic context or sig­
nalling state of a cell can affect the outcome of aberrant
ubiquitylation: NEDD4 functions as a tumour suppres­
sor in normal cells but as an oncogene in cells expressing
hyperactivated and/or mutant RAS.
Biochemical and functional analysis suggested
that the endosomal E3 RAB5 GDP/GTP exchange
factor (RABEX5) promotes monoubiquitylation and
diubiquity­lation of HRAS and NRAS, which induces
their localization to and retention in endosomes and
thus suppresses their signalling output73. In non-small-
cell lung cancer (NSCLC) cells expressing WT KRAS,
the DUB OTU domain-containing ubiquitin aldehyde-­
binding protein 1 (OTUB1) suppresses RAS ubiquity­
lation in a proteolysis-independent manner, leading to
increased RAS signalling74. However, in vitro studies
suggest that monoubiquitylation of KRAS at K147 or
HRAS at K117 can increase RAS signalling by inhibit­
ing GTPase-activating protein (GAP)-mediated hydro­
lysis or increasing GTP–GDP exchange, respectively75–77.
The expression of a KRAS‑G12V;K147L‑double mutant
protein in mouse fibroblast NIH‑3T3 cells reduces trans­
formation monitored in subcutaneous mouse models,
as compared to mice with KRAS‑G12V alone, indi­
cating that monoubiquitylation of KRAS‑G12V may
promote tumorigenesis75. However, as structural and
mutational studies were limited to human embryonic
kidney (HEK293) cells, further validation in cancer cells
is required, and the cellular contexts and the E3s and
DUBs that control the outcomes of monoubiquitylation
remain to be defined (FIG. 4b).
BRAF, the downstream effector kinase of RAS, is
also regulated by ubiquitylation. WT BRAF is targeted
for degradation by the E3 RNF149 (REF. 78). Ubiquitin-
dependent negative feedback control of BRAF activity
has been observed in Caenorhabditis elegans79, and
a similar mechanism has been described in human
cells80. ERK-mediated phosphorylation of either WT
BRAF or BRAF‑V600E, the most common mutant in
melanoma, primes it for ubiquitin-dependent degrada­
tion80. In C. elegans, this was shown to be mediated by
an FBXW7 homologue79, yet the human E3 remains to
be defined80. BRAF‑V600E, but not WT BRAF, requires
the chaperone heat shock protein 90 (HSP90) for proper
folding and stabilization81. When HSP90 activity is
inhibited pharmacologically, its client proteins, includ­
ing BRAF‑V600E, are degraded by cullin‑5-­mediated
ubiquitylation and proteasomal degradation81,82. CDH1
has also been suggested to limit BRAF signalling by
multiple mechanisms83. In non-malignant cells, the
APC/C–CDH1 complex induces BRAF proteolysis in
a cell-cycle-dependent manner, but this process is sup­
pressed in BRAF‑V600E‑expressing melanoma cells83.
Interestingly, ERK and CDK4 both phosphorylate CDH1,
REVIEWS
NATURE REVIEWS | CANCER
 VOLUME 18 | FEBRUARY 2018 | 77
©
 2
0
1
7
 M
a
c
m
i
l
l
a
n
 P
u
b
l
i
s
h
e
r
s
 L
i
m
i
t
e
d
,
 p
a
r
t
 o
f
 S
p
r
i
n
g
e
r
 N
a
t
u
r
e
.
 A
l
l
 r
i
g
h
t
s
 r
e
s
e
r
v
e
d
.
which decreases its association with APC/C and leads
to the accumulation of APC/C substrates, representing
a positive feedback loop between hyperactivated ERK
signalling and BRAF stability, as demonstrated in multi­
ple melanoma cell lines and in BRAF‑V600E‑expressing
immortalized mouse melanocytes. Consistent with this,
treatment of cells with ERK inhibitors or CDK inhibi­
tors decreases BRAF protein levels, suggesting that com­
bined inhibition of these kinases is a useful therapy for
­cancers with hyperactivated BRAF83. Importantly, CDH1
can also limit signalling in cancer cells expressing WT
BRAF (or dimerization-dependent BRAF mutants);
in this case, direct binding of CDH1 prevents BRAF
dimerization and full activation83. The dual-­suppressor
activity of CDH1 on BRAF illustrates the need to
develop CDC20‑specific inhibitors21,83. Although we
are yet to fully define ubiquitin-mediated regulation of
MAPK signalling, these examples provide the rationale
for developing novel therapies that limit RAS signal­
ling (for example, via OTUB1 inhibition) and suggest
a MAPK and PI3K–AKT–mTOR signalling
b Modulation of RAS by E3s
Effects of RAS monoubiquitylation
d E3s increase AKT activity
Proteasomal degradation
• Transcription and
   cell growth
• Cell cycle and
   DNA repair
• Translation and
    cell growth
• Autophagy
• Survival
• Glucose
   metabolism
e E3s modulate mTOR signalling
Proteasomal degradation
Regulation of mTORC formation
Regulation of protein activity
RAS–GTP
RAS
RAS
RAS
Ub
P
P
Ac
PI3K
PIP2
PIP3
AKT
RAS–GDP
RAF
RTK
MEK
ERK
PTEN
MDM2
PTPL1
mTOR
mTORC1
CKIα
GβL
GβL
GβL
RAPTOR
K63
K63
K63
PRAS40
mSIN1
mSIN1
mSIN1
RABEX5
Parkin
β-TRCP
DEPTOR
mTORC2
GβL
PROTOR
DEPTOR
DEPTOR
RICTOR
SCF–SKP2
SCF–SKP2
?
NF-κB
GSK3β
OTUB1
OTUD7B
TSC1 or
TSC2
mTOR
mTORC2
FOXO1 or
FOXO3
RAS
Increased GTP–RAS;
increased PI3K and
RAF binding
c Modulation of PI3K by E3s
Increase in PI3K signalling output
p85β
p85β
FBXL2
TRAF6
TRAF2
mTORC1 or
mTORC2
EGFR
PTEN
AKT
AKT
p300
SNO
REVIEWS
78 | FEBRUARY 2018 | VOLUME 18
www.nature.com/nrc
©
 2
0
1
7
 M
a
c
m
i
l
l
a
n
 P
u
b
l
i
s
h
e
r
s
 L
i
m
i
t
e
d
,
 p
a
r
t
 o
f
 S
p
r
i
n
g
e
r
 N
a
t
u
r
e
.
 A
l
l
 r
i
g
h
t
s
 r
e
s
e
r
v
e
d
.
possible mechanisms of action of therapeutics already
in clinical evaluation, including CDK4 and CDK6
­inhibitors and HSP90 inhibitors.
Regulation of PI3K–AKT–mTOR signalling. In contrast
to the relative paucity of studies on the regulation of the
MAPK pathway by ubiquitylation, multiple lines of evi­
dence support a role for ubiquitin-mediated degradative
and non-degradative pathways in regulating PI3K–
AKT–mTOR signalling (FIG. 4c,d,e). The SCF–F‑box
and leucine-rich repeat protein 2 (FBXL2) complex
has been demonstrated to regulate PI3K activation via
degradation of the regulatory subunit p85β84 (FIG. 4c).
Mechanistically, SCF–FBXL2 marks p85β for degrada­
tion following dephosphorylation of p85β by protein
tyrosine phosphatase PTPL1 (also known as PTPN13
and FAP1), thus preventing the competition between
free p85β and active PI3K (made up of p110 (the cata­
lytic subunit) and p85 heterodimers) for substrate bind­
ing84. Therefore, therapeutic targeting of FBXL2 may
represent one route to limit PI3K signalling in cancer84.
In addition to its effects on cyclins, parkin limits
AKT activity and WNT–β-catenin signalling (by tar­
geting β-catenin for proteasomal degradation), two
major pathways regulating cell growth and survival85.
Furthermore, PARK2 copy number loss has been
implicated in stimulating the PI3K–AKT pathway via
mitochondrial dysfunction in PTEN-expressing but not
PTEN-null cancers86: mechanistically, knockdown of
PARK2 in ­cancer cells impaired mitochondrial metabo­
lism, reflected by decreased ATP levels, increased oxi­
dative stress, activation of AMP-activated protein kinase
(AMPK) and phosphorylation and activation of endo­
thelial nitric oxide synthase (eNOS)86. The latter induces
PTEN S‑nitrosylation and concomitant ubiquitylation-­
dependent degradation86 (FIG. 4c). In agreement with
these in vitro findings86, both the Cancer Cell Line
Encyclopedia and The Cancer Genome Atlas (TCGA)
indicate that heterozygous deletion of PARK2 frequently
occurs in PTEN heterozygous cancer cells and tissues.
Accordingly, Pten+/– mice crossed with mice carrying a
targeted knockout of Park2 exon 3 are significantly more
tumour prone than Pten+/–Park2+/+ mice86.
The SCF–SKP2‑mediated K63‑linked ubiquityla­
tion of AKT1 and AKT2 promotes AKT recruitment
to ERBB receptors and its concomitant activation87.
ERBB activation in breast cancer cells elevates glucose
uptake and glycolysis and increases cellular resistance
to HER2‑targeting therapies87. Interestingly, these AKT-
mediated metabolic effects can be reversed by a SKP2
inhibitor, resulting in a p53‑independent but p27KIP1‑­
dependent senescence88. This inhibitor reduces the
viability of multiple cancer cell types (including p53‑­
deficient cancer cells) in vitro and in xenograft tumour
mouse models, pointing to a potential anti­tumour
activity of this inhibitor in SKP2‑overexpressing human
­cancers88. Of note, AKT1 can regulate SKP2 localization
by direct phosphorylation of SKP2 at S72 (REFS 89,90)
and via histone acetyltransferase p300 activation, which
in turn mediates SKP2 acetylation, thereby inhibiting
its interaction with CDH1 and enabling its retention
in the cytoplasm; this is associated with more aggres­
sive pheno­types in breast and prostate cancer91. Liver-
specific loss of HIPPO signalling in a liver cancer mouse
model also resulted in AKT–p300-mediated acetylation
and cytoplasmic retention of SKP2 (REF. 92). In hepato­
cytes isolated from mouse strains with defective HIPPO
signalling, cytoplasmic localization of SKP2 is associated
with p27KIP1 stabilization and, consequently, mitotic
arrest and polyploidy92. The effect of AKT activity on
SKP2 is also linked to the degradation of the transcrip­
tional activators forkhead box protein O1 (FOXO1) and
FOXO3, which positively regulate apoptosis; in turn,
this enables increased proliferation of polyploid cells,
genomic instability and increased formation of HCC in
the HIPPO signalling-deficient liver92.
The activity and stability of mTOR is regulated by a
diverse set of E3s (FIG. 4e): among them is SCF–FBXW7,
which binds, ubiquitylates and promotes mTOR degra­
dation93. In a subset of breast cancer cell lines, loss of a
single copy of FBXW7 appears mutually exclusive with
loss of a single copy of PTEN (a well-established, indi­
rect, negative regulator of mTOR), substantiating the
importance of FBXW7‑mediated mTOR stabilization in
tumorigenesis; thus, in addition to its effects on the cell
cycle, SCF–FBXW7 loss may potently stimulate anabolic
processes to promote tumour progression93.
Figure 4 | Regulation of mitotic signalling by ubiquitin ligases. a | Simplified
schematics of the PI3K–AKT–mTOR and MAPK pathways are depicted.  Ubiquitin ligases
(E3s) and deubiquitylating enzymes (DUBs) regulate degradation, activity or complex
assembly of receptor tyrosine kinases (RTKs), kinases and signalling molecules.
b | Monoubiquitylation or diubiquitylation modulates RAS activity. Elevated expression
of the DUB OTU domain-containing ubiquitin aldehyde-binding protein 1 (OTUB1) may
contribute to sustained RAS activation, as demonstrated in lung cancer74. c | High levels
of the PI3K regulatory subunit p85β compete with active PI3K (made up of p85–p110
heterodimers) for substrate binding, thereby limiting PI3K activity. Following
dephosphorylation of p85β by PTPL1, p85β is ubiquitylated by F‑box and leucine-rich
repeat protein 2 (FBXL2), leading to its proteasomal degradation and thereby increasing
PI3K signalling output84. d | PARK2 (which encodes E3 parkin) is frequently inactivated or
downregulated in cancer by diverse mechanisms. Parkin loss indirectly increases AKT
activity by affecting AKT upstream regulators. First, parkin loss leads to accumulation of
parkin substrates, including the epidermal growth factor receptor (EGFR). Second, parkin
loss impairs mitochondrial metabolism that results, from a cascade of events (indicated
by the dotted line), in PTEN S‑nitrosylation (by addition of the functional group
S‑nitrosothiol (SNO)), which primes PTEN for degradation. K63‑linked polyubiquitylation
of AKT by S-phase kinase-associated protein 1 (SKP1)–cullin 1–F‑box protein (SCF)–SKP2
directly increases its activity. AKT can regulate stabilization and localization of SKP2,
wherein AKT either directly phosphorylates SKP2 (REFS 89,90) (not shown) or mediates
histone acetyltransferase p300 activation, which in turn leads to SKP2 acetylation (Ac)
and thus stabilization and cytoplasmic retention of SKP2 (REF. 91). e | A casein kinase Iα
(CKIα)–SCF–β-transducin repeat-containing protein (β-TRCP) auto-amplification loop
mediates full activation of mTOR via DEP domain-containing mTOR-interacting protein
(DEPTOR) degradation199. The dynamic assembly of mTOR complex 1 (mTORC1) and
mTORC2 is under the control of ubiquitylation. Ubiquitin-dependent deregulation of
these signalling hubs in turn alters cellular processes as diverse as cell growth,
metabolism, DNA repair, transcription, translation and survival. FOXO, forkhead box
protein O; GSK3β, glycogen synthase kinase 3β; NF‑κB, nuclear factor-κB; OTUD7B, OTU
domain-containing protein 7B; P, phosphorylation; PIP, phosphatidylinositol phosphate;
PRAS40, proline-rich AKT1 substrate 1; PROTOR, proline-rich protein; RABEX5, RAB5
GDP/GTP exchange factor; RAPTOR, regulatory-associated protein of mTOR; RICTOR,
rapamycin-insensitive companion of mTOR; TRAF, tumour necrosis factor (TNF)-
receptor-associated factor; TSC1, tuberous sclerosis 1 (also known as harmartin);
TSC2, also known as tuberin; Ub, ubiquitin.
◀
REVIEWS
NATURE REVIEWS | CANCER
 VOLUME 18 | FEBRUARY 2018 | 79
©
 2
0
1
7
 M
a
c
m
i
l
l
a
n
 P
u
b
l
i
s
h
e
r
s
 L
i
m
i
t
e
d
,
 p
a
r
t
 o
f
 S
p
r
i
n
g
e
r
 N
a
t
u
r
e
.
 A
l
l
 r
i
g
h
t
s
 r
e
s
e
r
v
e
d
.
Hypoxic tension
The level of oxygen (usually
measured as a percentage)
in a given tissue or
microenvironment. The lower
the level of oxygen, the
higher the tension.
Unfolded protein response
(UPR). A well-defined process
that plays a critical role in
restoring homeostasis following
accumulation of potentially
toxic misfolded proteins in the
endoplasmic reticulum.
Succination
A process wherein fumarate
reacts with cysteine residues in
proteins by a Michael addition
reaction to form S-(2‑succinyl)
cysteine.
Activation of mTOR complex 1 (mTORC1) was
previously shown to be regulated by K63‑linked
ubiquity­lation mediated by the E3 tumour necrosis
factor (TNF)-receptor-associated factor 6 (TRAF6) in
complex with the scaffold protein p62 (also known as
sequestosome 1)94, and more recent work has demon­
strated the role of ubiquitylation in the dynamic
assembly of mTORC1 and mTORC2 (REF. 95). TRAF2
mediates K63‑linked ubiquitylation of GβL (also known
as mLST8; a component of both mTORCs) and disrupts
its binding to the mTORC2‑specific component mSIN1
(also known as MAPKAP1), thereby restrictin
\end{document}
